\documentclass[../main]{subfiles}
\begin{document}
\ifSubfilesClassLoaded{\setcounter{page}{40}}{}
\section{Комунізм і управління суспільним життям}
\label{sec:07_kommunizm_upravlenie}

Виробництво та розподіл вільного часу в перехідний період до комунізму потребують державного регулювання. Однак держава потрібна лише до тих пір, поки виробництво вільного часу збігається з виробництвом і розподілом матеріальних благ, необхідних людям, і до тих пір, поки саме це виробництво матеріальних благ ще не здійснюється на комуністичних засадах. Сам же вільний час, виступаючи, власне, як вільний (час комунізму), виходить за межі будь-якого зовнішнього, відносно нього, державного управління.

У сфері вільного часу держава не існує навіть зараз, за часів капіталізму. Там, де відбувається вільна комуністична діяльність, управління здійснюється всіма членами вільної асоціації людей. У міру збільшення частки комуністичної діяльності держава відмирає.

Сприяння розвитку вільного часу та безпосередньо колективних форм вільної діяльності є актом знищення умов для існування класів у перехідний період, і в тій самій мірі актом самознищення держави. Держава – це лише апарат забезпечення інтересів панівного класу, який здійснює управління в інтересах цього класу, а ця його основна функція відпадає в міру знищення класів, зокрема й панівного.

Таким чином, завдання держави в перехідний період зводиться до забезпечення обмеження (знищення) товарного характеру виробництва та піклування про створення умов для комуністичної вільної діяльності членів суспільства.

Перша з них виникає і перед капіталістичними державами. Капіталізм постійно породжує війни, а логіка війни вимагає системи прямого розподілу продуктів, призначених для особистого споживання. Держава, яка не справляється з цим завданням, приречена на знищення (або як держава, або разом із населенням). Завдання налагодження такої системи передбачає в ряді пунктів обмеження окремих капіталістів в інтересах капіталістів як класу, що не завжди може бути здійснено в повній мірі та з достатньою ефективністю.

Під час війни особисте забезпечення – як продуктами харчування, так і речами першої необхідності – не може здійснюватися на засадах капіталістичного приватного присвоєння продуктів праці. Солдат, навіть якщо він найманець, під час війни має бути забезпечений своїм роботодавцем, принаймні, зброєю, провізією та обмундируванням. В іншому випадку армія перетворюється на бандитські угруповання. Але ці угруповання не здатні вести сучасну війну. Для ведення війни сучасній армії потрібна високотехнологічна зброя, предмети особистого вжитку, необхідні для відновлення життєвих сил солдата, а також засоби особистого захисту. І навіть якщо все це забезпечується шляхом безпосереднього грабунку населення окупованих територій, це не може бути особистою справою окремих солдатів. Це лише умова, необхідна для того, щоб він міг виконувати свою роботу. Але ця умова відрізняється від умов найманої праці, що ґрунтується на приватному характері присвоєння. Присвоєння продуктів праці в усіх їхніх істотних і значущих для соціального відтворення індивіда аспектах у класичних капіталістичних відносинах купівлі-продажу робочої сили виноситься за межі виробництва, що ґрунтується на найманій праці. Особисте споживання визначається приватним характером присвоєння продуктів праці, опосередкованим грошима. І це, у свою чергу, гарантує, що працівник знову і знову буде змушений приступати до праці, щоб забезпечити себе всім необхідним. На війні все відбувається зовсім інакше.

Споживання в цьому випадку не може бути навіть формально залишено на розсуд самих індивідів. Отже, сучасна війна – це, окрім усього іншого, конкуренція форм нетоварного виробництва, що склалися ще за часів капіталізму, – безпосереднього виробництва людиною. І тут мається на увазі як армія як власне військова сила, так і трудова армія, яка формується по-різному. Акцент зміщується. Якщо для нормального перебігу капіталізму на перший план виходить те, що людина повинна робити свою справу, щоб жити, то тут, навпаки, людина повинна жити, щоб робити свою справу. Держава змушена брати на себе відповідальність за це.

При цьому грошово-товарні відносини стають очевидною перешкодою для виконання цих функцій. Війна робить неможливими старі, довоєнні способи та форми діяльності людей. У таких умовах держава об’єктивно посилюється, бере на себе функції розподілу в процесі безпосереднього виробництва, як в армії, так і в тилу, займається налагодженням і підтримкою господарських зв’язків. Але держава як апарат придушення одного класу іншим, точніше, придушення меншістю більшості, в принципі не може вирішити ці завдання в повній мірі, і тим більше не ставить перед собою ті завдання, які не слугують інтересам правлячого класу. Але саме такого роду завдань (які розходяться з інтересами правлячого класу) під час війни у суспільства стає дуже багато. Вони об’єктивно постають перед людьми, тому що від їх вирішення залежить саме їхнє фізичне існування. Держава часто навіть не намагається вирішувати ці проблеми, залишаючи їх на розсуд різноманітних недержавних організацій та об’єднань людей. Це змушує мільйони людей шукати нові форми організації, які б брали на себе державні функції. І ці органи громадського самоврядування, які беруть на себе функції держави, у широкому сенсі, є і не є державою.

Найважливіша і визначальна з проблем, вирішення якої вимагає нових організаційних форм на противагу капіталістичній державі, – це запобігання та припинення війни, а також мінімізація для суспільства втрат і руйнувань, спричинених нею. Реалізація цього найважливішого завдання якраз і суперечить інтересам капіталу, який ще не отримав з війни всі ті вигоди, які міг би отримати. Підкреслимо, що тут важливим може бути не війна сама по собі, не насильство саме по собі, а опір насильству реакційного класу, який здійснюється угрупованнями цього класу в ім’я своїх корисливих інтересів, ще й руками іншого класу. В опорі та формах опору (хоча й не лише в них), а не у війні самій по собі, криється найважливіша тенденція до самоорганізації на підступах до революції. Без такої самоорганізації, розроблені в рамках капіталізму способи безпосереднього, не опосередкованого матеріальними цінностями виробництва людини та економічні механізми їх реалізації, не можуть стати відправною точкою переходу суспільства на нові засади розвитку.

Ці форми самоорганізації, породжені тенденцією до війни та самою війною (як протистояння війні), що використовуються для управління суспільством, є перехідними і можуть розвиватися у трьох напрямках: 1) інтеграція з державою; 2) деградація та зникнення (інтеграція не завжди означає деградацію); 3) перейняття на себе всіх функцій держави, перетворення на нову державну владу. Таким чином, ці організаційні форми об’єктивно є формами класової боротьби там, де вони змушені протистояти державі, переймаючи на себе функції тих чи інших державних органів, а отже, позбавляючи апарат політичного панування правлячого класу повноважень. І якщо в певних умовах це фактично полегшення для цих органів, оскільки самі вони не можуть виконувати свої функції, то подальший рух цим шляхом неминуче наштовхується на протидію держави, на спроби обмежити функції та вплив форм самоорганізації населення. Яким буде результат і характер цього протистояння, залежить від конкретних обставин боротьби та дій сторін, що борються.

Якщо розглядати історію XX століття, то однією з найефективніших форм самоорганізації, які брали на себе функції держави, були ради. Ради робітників, солдатів і селян одночасно були представницько-законодавчими, розпорядчими та контрольними органами. І хоча ради виникли ще в 1905 році, їхнє реальне перетворення на державну силу, двовладдя, і тим більше гасло «Вся влада радам», а також конкретні кроки від гасла до його реалізації могли виникнути лише в умовах політичного банкрутства старої держави в процесі виснажливої війни. Після лютневої революції банкрутство держави як апарату, який претендував на те, щоб організовувати життя суспільства, стало очевидним, оскільки не лише царський уряд, а й Тимчасовий уряд виявився безсилим у цьому відношенні. Водночас почали формуватися та активно залучатися до справ як місцевого самоврядування, так і державного управління ради. Створювалася ситуація двовладдя як у тилу, так і в армії. У зв’язку з цим особливо важливо враховувати можливість у майбутньому повторення ситуації, яка склалася навесні та на початку літа 1917 року. Тоді був можливий мирний шлях передачі влади радам як крок на шляху до соціальної революції — створення умов для соціалізму. І ця можливість існувала завдяки двом обставинам: силі та суспільній значущості самих рад, а також озброєнню населення в результаті війни. Лозунг «Вся влада радам» у той час був констатацією наявності такого мирного шляху революції. І лише після липневих подій 1917 року мирний шлях став неможливим. Але сам факт такої можливості для нас є найважливішим. Адже якщо складуться такі умови, дуже важливо їх не проґавити, тим паче, враховуючи розвиток сучасної зброї, мирний шлях цілком може виявитися єдино можливим для революції.

У сучасних умовах можуть виникати й виникають інші форми, не завжди пов’язані з війною, але завжди з необхідністю вирішувати важливі суспільні проблеми. Суть цих перехідних форм скрізь однакова: населення безпосередньо бере на себе функції державної влади, тим самим скасовуючи її у звичному для капіталізму розумінні.

Тепер перейдемо до розгляду управління суспільним життям за комунізму. Управління суспільним життям взагалі зводиться до планування та коригування плану з урахуванням реальних змін, вираження волі (розпоряджень), обліку та контролю над керованим процесом, безпосереднього втручання в разі необхідності.

За часів комунізму планування передбачало планування розвитку людини, яка ставала головною метою суспільного виробництва, планування створення нових відносин між людьми. Таке планування збігається з прогнозуванням у тому сенсі, що найнадійніший прогноз – це той, для реалізації якого докладаються цілеспрямовані зусилля. Усе це суспільство має навчитися робити без нагляду з боку органу «управління», який насправді є інструментом пригнічення, що називається державою.

У міру становлення нової колективності всі її функції, пов’язані з виробництвом матеріальних благ, мають бути передані цим благам. Держава в перехідний до комунізму період не просто відмирає, вона відмирає лише в тому випадку, якщо працює на власне відмирання, в тій мірі, в якій вона піклується про те, щоб люди обходилися без неї. Усі залишки неволі, в тому числі й способи використання часу, з якими пов’язане існування держави, протистоять вільному часу як простору людського розвитку, протистоять вільній людській діяльності. Там, де можлива повна свобода, немає місця для держави. Але там, де повна свобода ще неможлива безпосередньо, наявність певного типу держави не тільки фіксує факт цієї неможливості, але й є інструментом для її подолання.

Проти заперечення необхідності існування держави висувається той самий аргумент, що й проти відмови від поділу праці. Суть його полягає в тому, що суспільне життя стає надзвичайно складним і постійно ускладнюється, а отже, для управління таким складним життям потрібна диференціація функцій і спеціальний апарат. Але з цієї складності зовсім не випливає необхідність у створенні особливих, над суспільством поставлених, загонів чиновників і озброєних людей. Якби не було розколу суспільства на класи, самокерована та самоорганізована асоціація людей була б можливою, попри всю її складність, яка відповідала б рівню складності суспільного життя, на чому, слідуючи за Марксом та Енгельсом, наполягав Ленін.

Тим більше, що, говорячи про комунізм, ми маємо на увазі суспільство, де подолано поділ праці, де люди більш-менш однаково розвинені щодо здоров’я, чуттєвості, мислення та моралі, мають політехнічну, науково-теоретичну, універсальну освіту на сучасному необхідному (тобто найвищому) для суспільства рівні, а отже, можуть за короткий час опановувати будь-яку окрему «професію» та застосовувати свої сили там, де це потрібно, і де це їм самим цікаво, хай навіть один здебільшого все життя присвячує вивченню птахів, а інший – балету.

І саме тому в управлінні суспільним розвитком із використанням досягнутих технологій, спираючись на сформовану людством культуру мислення, за комунізму беруть участь усі. Це найважливіша умова універсального розвитку всіх членів суспільства – протягом усього життя всі повинні вміти керувати та мати звичку керувати суспільними справами – усіма справами, в яких самі беруть участь. Таке суспільство більш реальне, ніж фантастичний класовий світ і «гармонійні взаємини» між класами, які пропагують сучасні буржуазні ідеологи. Навпаки, марксизм наполягає на тому, що сама наявність держави говорить про те, що цей світ в умовах існування класів неможливий. Для поступового зникнення держави повинні існувати відповідні соціальні та технічні передумови.

Що стосується технічних передумов, то за умови наявності якісної загальнообов’язкової політехнічної освіти, що базується на сучасних технологіях, організація управління суспільними процесами не була б надто складною в суспільстві, де відсутня приватна власність (тому найважливішою передумовою є формальне об’єднання засобів виробництва). Звісно, для цього необхідно усунути соціальні бар’єри, що також вимагає від суспільства серйозної роботи і не є чимось, що можна досягти одразу. Головне, що чисто технічно вже на цьому етапі це не видається неможливим. Уся інформація про стан виробництва та розподілу в усіх сферах суспільного життя в режимі реального часу, з можливістю відстеження ланцюжків до самого початку, з урахуванням ресурсів, потреб, технологічних можливостей, можливостей зв’язку та голосування тощо. від найнижчого рівня (безпосередньої участі) до найвищого рівня колективного прийняття рішень на рівні як локальних спільнот, так і всього суспільства – все це в режимі реального часу може бути доступне кожному члену асоціації.

Голосування та інші, звичні нам прояви демократії потрібні лише на початковому етапі. У процесі переходу від управління людьми до управління речами та процесами потреба в голосуваннях поступово відпадає, поступаючись місцем об’єктивній необхідності, що міститься в самих речах і в їх взаємодії. На це вказував Ленін у праці «Держава і революція», детально аналізуючи спадщину Маркса та Енгельса: «Про «відмирання» і навіть ще яскравіше та виразніше – про «засинання» Енгельс говорить цілком чітко та однозначно щодо епохи після «передачі засобів виробництва у власність держави від імені всього суспільства», тобто після соціалістичної революції. Ми всі знаємо, що політичною формою «держави» в цей час є найповніша демократія. Але жодному з опортуністів, які безсоромно спотворюють марксизм, не спадає на думку, що тут, отже, в роботах Енгельса, йдеться про «засинання» і «відмирання» демократії. На перший погляд це здається дуже дивним. Але це «незрозуміло» лише для того, хто не замислювався над тим, що демократія також є державою і, отже, демократія теж зникне, коли зникне держава. Буржуазну державу може «знищити» лише революція. Держава взагалі, тобто найбільш повна демократія, може лише «відмерти».

У сучасну епоху бурхливого розвитку технологій організація поступового зникнення держави не здається вже такою складною справою з суто технічної точки зору. Облік і контроль у сфері управління речами та процесами за таких умов стають питанням техніки, а їхні результати стають доступними кожному. У цьому процесі відпадає потреба в ієрархії та зникає передумова для того, щоб ці функції могли відокремитися та перетворитися на окрему професію. За умови впровадження автоматизованих систем для управління та забезпечення повного доступу до баз даних для кожного та можливості працювати з ними, зникає потреба в «особливій армії чиновників», які могли б навіть «зовнішньо» стояти «над» суспільством як окремий «апарат».

Таким чином буде вирішено одне з найважливіших завдань переходу до соціалізму – налагодження обліку та контролю, на важливості якого постійно наголошував Ленін: «Облік і контроль – ось головне, що потрібне для «налагодження», для правильного функціонування першої фази комуністичного суспільства». Йдеться про такий облік і контроль над суспільним виробництвом, організація якого вже сама по собі означає організацію нових форм колективності. «Якщо справді всі беруть участь в управлінні державою, тут вже капіталізму не утриматися. І розвиток капіталізму, у свою чергу, створює передумови для того, щоб справді всі могли брати участь в управлінні державою». До таких передумов належить повсюдне поширення грамотності, здійснене вже в низці найбільш передових капіталістичних країн, а також «навчання та дисциплінування» мільйонів робітників за допомогою великих, складних, об’єднаних механізмів поштової, залізничної, великої фабричної та торгової систем, банківської справи тощо.

Ми можемо ще раз наголосити, що до таких передумов належить вже зараз розвиток комп’ютерної техніки та закладені в ній можливості, які в принципі не можуть бути використані за капіталізму, коли немає єдиного суспільного господарства, а зв’язок між виробниками здійснюється через ринок. Ці можливості в організації управління суспільним життям досліджував у своїх розробках Віктор Михайлович Глушков, який розробляв Загальнодержавну систему управління економікою в СРСР. Глушков розглядав управління суспільним виробництвом в цілому, за допомогою систем управління, де кожен має доступ до найдетальнішої картини виробництва в реальному часі, а там, де людські рішення не потрібні, система діє автоматично за заздалегідь визначеними людьми схемами відповідно до суспільних цілей.

Чим більше розвиватимуться технології, тим простіше буде організувати управління та забезпечити участь у ньому всіх членів суспільства. Щодо змістовної сторони, то вона вимагає відповідного розвитку членів суспільства. Складність полягає в тому, що рух до комунізму потрібно починати не з людьми комуністичного завтра, а з людьми, які сформувалися в умовах розподілу праці, і які долають свою обмеженість за допомогою освіти та колективної діяльності.

Політехнічна освіта стає необхідністю для членів вільної асоціації. Ідея політехнізму досить добре розроблена, тож ми можемо сміливо порекомендувати читачеві відповідну літературу. Тут же зазначимо лише те, що для управління суспільними процесами члени суспільства повинні володіти цілісною науковою картиною світу та універсальними навичками застосування своїх знань. Сучасному чиновнику для виконання своєї функції наукова картина світу навіть у першому наближенні абсолютно не потрібна, тому вона у нього сама собою зникає, як щось зайве, навіть якщо колись була. Чиновник – лише жива часткова деталь застарілої державної машини. «Суспільство, яке по-новому організує виробництво на основі вільної та рівної асоціації виробників, відправить всю державну машину туди, де їй буде тоді належне місце: у музей старожитностей, поруч із прялкою та бронзовим сокиром», – говорив Енгельс.

Звісно, сама задача управління суспільними процесами в такій асоціації та свідомого побудови суспільного життя вимагатиме від людей дуже багато, зокрема розвиненого діалектичного мислення, розвиненої чутливості, моральності та відповідного рівня грамотності. Але це не недосяжні висоти, доступні лише для обраних, а те, що за відповідних умов стане надбанням кожного. У міру залучення людей до активного суспільного життя як суб’єктів цього життя, вони виходитимуть на передній план людської культури. Так ми розуміємо відомий афоризм Леніна про те, що кожна кухарка має навчитися управляти державою. Перешкодою для цього є лише спосіб діяльності сучасних людей, який перетворює одну людину на державного чиновника, а іншу – на кухарку.

Комуністичне суспільство – це вільне об’єднання людей, яке самостійно керує своїм розвитком за розумним планом.

Яким саме чином воно буде керувати собою, за допомогою яких саме засобів та організаційних заходів? На це питання ми відповісти не можемо. Це – справа цього об’єднання, яке буде змушене, виходячи з необхідності, розвивати відповідні організаційні форми в конкретних історичних умовах. Самоуправління розвиватиметься разом із розвитком комуністичного суспільства.
\end{document}
